\documentclass[11pt]{article}

\usepackage[margin=1in]{geometry}
\usepackage{amsmath,amssymb,bm}
\usepackage{booktabs}
\usepackage{array}
\usepackage{setspace}

\setstretch{1.08}

\title{MPC Formulation and Current Parameters}
\author{}
\date{}

\begin{document}

\maketitle

\section{Objective Function}

The current linear time-varying MPC problem is written as

\begin{equation}
\min_{\{X_k,U_k\}_{k=1}^{N}} J
=
J_{\mathrm{ref}}
+
J_{\mathrm{lane}}
+
J_{\mathrm{repulsive}}
+
J_{\mathrm{control}}
\end{equation}

where the state and control are

\begin{equation}
X_k = [x_k,\; y_k,\; v_k,\; \psi_k]
\end{equation}

\begin{equation}
U_k = [a_k,\; \delta_k].
\end{equation}

\subsection{Attractive Reference Cost}

The destination-tracking term is

\begin{equation}
J_{\mathrm{ref}}
=
w_{\mathrm{att}}
\sum_{k=1}^{N}
\left[
q_x (x_k-x_{\mathrm{ref}})^2
+
q_y (y_k-y_{\mathrm{ref}})^2
+
q_v (v_k-v_{\mathrm{ref}})^2
+
q_{\psi}\big(\psi_k-\psi_{\mathrm{ref}}\big)^2
\right].
\end{equation}

\subsection{Lane-Center Following Cost}

The lane-center lateral error is

\begin{equation}
e_{y,k}
=
-
\big(x_k-x^{\mathrm{lane}}_{\mathrm{ref},k}\big)
\sin\theta^{\mathrm{lane}}_{\mathrm{ref},k}
+
\big(y_k-y^{\mathrm{lane}}_{\mathrm{ref},k}\big)
\cos\theta^{\mathrm{lane}}_{\mathrm{ref},k}.
\end{equation}

The stage-varying lane weight is

\begin{equation}
w_k = w_0 \alpha^{k-1}.
\end{equation}

The lane term is

\begin{equation}
J_{\mathrm{lane}}
=
\sum_{k=1}^{N}
w_k
\left[
e_{y,k}^2
+
q_{\psi,\mathrm{lane}}
\big(\psi_k-\psi^{\mathrm{lane}}_{\mathrm{ref},k}\big)^2
\right].
\end{equation}

\subsection{Control Smoothness Cost}

The control-rate smoothing term is

\begin{equation}
J_{\mathrm{control}}
=
w_{\mathrm{ctrl}}
\sum_{k=1}^{N}
\left[
q_a\left(\frac{a_k-a_{k-1}}{\Delta t}\right)^2
+
q_{\delta}\left(\frac{\delta_k-\delta_{k-1}}{\Delta t}\right)^2
\right].
\end{equation}

\section{Current Repulsive Potential Field}

The live obstacle cost is no longer the old Gaussian potential. It is now a
super-ellipsoid-based potential field with separate safe-zone and collision-zone
costs.

\subsection{Relative Pose in the Obstacle Frame}

For obstacle center $(x_o,y_o)$ and obstacle heading $\psi_o$, the ego position
is transformed into the obstacle-local frame as

\begin{equation}
\begin{bmatrix}
x_{\mathrm{loc}} \\
y_{\mathrm{loc}}
\end{bmatrix}
=
R(-\psi_o)
\begin{bmatrix}
x_e-x_o \\
y_e-y_o
\end{bmatrix}.
\end{equation}

\subsection{Projected Ego Shape}

Let the relative heading be

\begin{equation}
\Delta\psi = \psi_e-\psi_o.
\end{equation}

The ego footprint is projected onto the obstacle frame:

\begin{equation}
L_e^{\mathrm{proj}}
=
\left|L_e\cos(\Delta\psi)\right|
+
\left|W_e\sin(\Delta\psi)\right|
\end{equation}

\begin{equation}
W_e^{\mathrm{proj}}
=
\left|L_e\sin(\Delta\psi)\right|
+
\left|W_e\cos(\Delta\psi)\right|.
\end{equation}

The static inflated base half-sizes are

\begin{equation}
x_0
=
\frac{L_e^{\mathrm{proj}} + L_o}{2}
+
b_{\mathrm{long}}
\end{equation}

\begin{equation}
y_0
=
\frac{W_e^{\mathrm{proj}} + W_o}{2}
+
b_{\mathrm{lat}}.
\end{equation}

\subsection{Dynamic Inflation}

The relative approach speeds are

\begin{equation}
\Delta u = \max\!\big(0,\; v_e\cos(\Delta\psi)-v_o\big)
\end{equation}

\begin{equation}
\Delta v = \max\!\big(v_{\mathrm{lat,min}},\; |v_e\sin(\Delta\psi)|\big).
\end{equation}

The collision-zone half-sizes are

\begin{equation}
x_c = x_0 + \frac{\Delta u^2}{2a_{\max}}
\end{equation}

\begin{equation}
y_c = y_0 + \frac{\Delta v^2}{2a_{\max}}.
\end{equation}

The safe-zone half-sizes are

\begin{equation}
x_s = x_0 + \Delta u T_r + \frac{\Delta u^2}{2a_{\mathrm{comfort}}}
\end{equation}

\begin{equation}
y_s = y_0 + \Delta v T_r + \frac{\Delta v^2}{2a_{\mathrm{comfort}}}.
\end{equation}

\subsection{Geometric Limits on the Field Size}

The current project also limits the field dimensions.

The longitudinal half-lengths are capped by a maximum full field length
$L_{\max}$:

\begin{equation}
x_c \leftarrow \min\left(x_c,\; \frac{L_{\max}}{2}\right),
\qquad
x_s \leftarrow \min\left(x_s,\; \frac{L_{\max}}{2}\right).
\end{equation}

The lateral half-widths are capped by the larger of the lane-width limit and
the obstacle width:

\begin{equation}
W_{\mathrm{lat,max}} = \max\!\big(W_{\mathrm{lane}} f_{\mathrm{lane}},\; W_o\big).
\end{equation}

Therefore,

\begin{equation}
y_c \leftarrow \min\left(y_c,\; \frac{W_{\mathrm{lat,max}}}{2}\right),
\qquad
y_s \leftarrow \min\left(y_s,\; \frac{W_{\mathrm{lat,max}}}{2}\right).
\end{equation}

\subsection{Normalized Super-Ellipsoid Distances}

With super-ellipsoid exponent $n$, the normalized collision and safe distances are

\begin{equation}
r_c
=
\left[
\left|\frac{x_{\mathrm{loc}}}{x_c}\right|^n
+
\left|\frac{y_{\mathrm{loc}}}{y_c}\right|^n
\right]^{1/n}
\end{equation}

\begin{equation}
r_s
=
\left[
\left|\frac{x_{\mathrm{loc}}}{x_s}\right|^n
+
\left|\frac{y_{\mathrm{loc}}}{y_s}\right|^n
\right]^{1/n}.
\end{equation}

\subsection{Stage Obstacle Cost}

For each obstacle $i$ at stage $k$, the current repulsive cost is

\begin{equation}
J_{\mathrm{safe}}^{(i,k)}
=
w_s \exp\!\left[-k_s\big(r_s^{(i,k)}-s_s\big)\right]
\end{equation}

\begin{equation}
J_{\mathrm{collision}}^{(i,k)}
=
w_c \exp\!\left[-k_c\big(r_c^{(i,k)}-s_c\big)\right]
\end{equation}

\begin{equation}
J_{\mathrm{obs}}^{(i,k)}
=
J_{\mathrm{safe}}^{(i,k)}
+
J_{\mathrm{collision}}^{(i,k)}.
\end{equation}

Therefore, the total repulsive term over the horizon is

\begin{equation}
J_{\mathrm{repulsive}}
=
\sum_{k=1}^{N}\sum_{i=1}^{M} J_{\mathrm{obs}}^{(i,k)}.
\end{equation}

\paragraph{Interpretation.}
Both the safe-zone and collision-zone costs are always active. As the ego gets
closer to the obstacle, the normalized distances $r_s$ and $r_c$ become smaller,
which increases the exponential cost. The parameters $k_s,k_c$ control how
quickly the cost changes with distance, while $s_s,s_c$ shift the effective
distance level at which the exponent becomes significant.

\section{Lookahead Distance and Rolling Temporary Destination}

\subsection{Current CARLA-Branch Flow}

The rolling temporary destination is updated in the main MPC loop, not in the
behavior-planner worker. The current flow is:

\begin{enumerate}
    \item Read the current ego state $(x_e, y_e, v_e, \psi_e)$.
    \item Rebuild the global route from the current ego position to the final
    destination when the route-rebuild timer expires.
    \item Compute the lookahead distance used for the rolling temporary
    destination.
    \item Select the target lane from the most recent accepted behavior:
    \begin{itemize}
        \item \texttt{LANE\_KEEP} or \texttt{GO\_STRAIGHT}: current ego lane
        \item \texttt{LANE\_CHANGE\_LEFT} or \texttt{TURN\_LEFT}: adjacent left lane
        \item \texttt{LANE\_CHANGE\_RIGHT} or \texttt{TURN\_RIGHT}: adjacent right lane
    \end{itemize}
    \item Starting from the ego's current position on that lane, walk forward
    along the lane-center waypoint chain by the lookahead distance.
    \item Interpolate along the final segment to place the temporary
    destination exactly at that forward distance.
\end{enumerate}

Thus, at every control step,

\begin{equation}
\mathbf{p}_{\mathrm{temp}}(k)
=
\operatorname{LanePointAhead}
\Big(
\mathbf{x}_{\mathrm{ego}}(k),\;
\ell_{\mathrm{selected}}(k),\;
d_{\mathrm{lookahead}}(k)
\Big),
\end{equation}

where $\ell_{\mathrm{selected}}(k)$ is the lane selected by the latest
behavior-planner decision.

\subsection{Lookahead-Distance Helper}

The code still contains the dynamic lookahead helper

\begin{equation}
d_{\mathrm{raw}} = d_{\min} + k_v v_e - k_c |\kappa_{\mathrm{route}}|,
\end{equation}

followed by clamping:

\begin{equation}
d_{\mathrm{lookahead}}
=
\operatorname{clip}\!\left(d_{\mathrm{raw}},\; d_{\min},\; d_{\max}\right).
\end{equation}

Here,
\begin{itemize}
    \item $d_{\min}$ is the minimum lookahead distance,
    \item $d_{\max}$ is the maximum lookahead distance,
    \item $k_v$ is the speed gain,
    \item $k_c$ is the curvature penalty gain,
    \item $\kappa_{\mathrm{route}}$ is estimated from future route samples.
\end{itemize}

\subsection{Route Curvature Calculation}

The route curvature used in the lookahead helper is estimated from three points
sampled along the current global route:

\begin{align}
\mathbf{p}_1 &= \mathbf{r}(s + \Delta s), \\
\mathbf{p}_2 &= \mathbf{r}(s + 2\Delta s), \\
\mathbf{p}_3 &= \mathbf{r}(s + 3\Delta s),
\end{align}

where
\begin{itemize}
    \item $s$ is the ego's current progress along the route,
    \item $\mathbf{r}(\cdot)$ is the route point sampled at a given progress,
    \item $\Delta s$ is the configured curvature sample spacing.
\end{itemize}

The side lengths of the triangle formed by these samples are

\begin{align}
a &= \|\mathbf{p}_2 - \mathbf{p}_1\|, \\
b &= \|\mathbf{p}_3 - \mathbf{p}_2\|, \\
c &= \|\mathbf{p}_3 - \mathbf{p}_1\|.
\end{align}

The triangle area is

\begin{equation}
A_{\triangle}
=
\frac{1}{2}
\left|
x_1(y_2-y_3) + x_2(y_3-y_1) + x_3(y_1-y_2)
\right|.
\end{equation}

The route curvature estimate is then

\begin{equation}
\kappa_{\mathrm{route}} = \frac{4A_{\triangle}}{abc}.
\end{equation}

If any of the side lengths are numerically too small, the code sets
\[
\kappa_{\mathrm{route}} = 0.
\]

Therefore, sharper route bending increases $|\kappa_{\mathrm{route}}|$ and
reduces the dynamic lookahead distance through the term
$-k_c |\kappa_{\mathrm{route}}|$.

\subsection{Current Override Used in the CARLA Test Scenarios}

For the current CARLA test scenarios, the main loop explicitly overrides the
dynamic helper with a scenario-level fixed distance:

\begin{equation}
d_{\mathrm{lookahead}} = d_{\mathrm{fixed}}.
\end{equation}

At present,

\begin{equation}
d_{\mathrm{fixed}} = 50~\mathrm{m}.
\end{equation}

Therefore, the temporary destination is rebuilt every main-loop step as the
point $50$ meters ahead of the ego vehicle on the currently selected lane.

\subsection{Route-Based Fallback}

If lane-based destination construction fails, the current runner falls back to
the route-based temporary destination sampled from the global route:

\begin{equation}
s_{\mathrm{target}}
=
\min\!\left(
s_{\mathrm{route,max}},
\;
s_{\mathrm{ego}} + d_{\mathrm{lookahead}}
\right),
\end{equation}

where $s_{\mathrm{ego}}$ is the current progress along the global route and
$s_{\mathrm{route,max}}$ is the total route length.

\subsection{Implication}

The intended behavior is:
\begin{itemize}
    \item the behavior planner decides only the lane action,
    \item the main MPC loop updates the temporary destination every step,
    \item and the blue dot should remain a rolling target ahead of the ego
    vehicle rather than a waypoint that updates only when reached.
\end{itemize}

\section{Current Parameter Values}

The current values below are taken from the live \texttt{MPC/mpc.yaml} file.

\subsection{Planner and Timing}

\begin{align}
\text{Horizon} &= 3.0~\mathrm{s} \\
\Delta t &= 0.05~\mathrm{s} \\
\text{Planning frequency} &= 2~\mathrm{Hz}
\end{align}

\subsection{Attractive Cost Parameters}

\begin{align}
w_{\mathrm{att}} &= 1.0 \\
[q_x,\; q_y,\; q_v,\; q_{\psi}] &= [10,\; 0,\; 1,\; 0]
\end{align}

\subsection{Lane-Center Cost Parameters}

\begin{align}
w_0 &= 30.0 \\
\alpha &= 0.97 \\
q_{\psi,\mathrm{lane}} &= 10.0
\end{align}

\subsection{Control Cost Parameters}

\begin{align}
w_{\mathrm{ctrl}} &= 30.0 \\
[q_a,\; q_{\delta}] &= [10,\; 10]
\end{align}

\subsection{Repulsive Potential Parameters}

\begin{align}
w_s &= 100 \\
w_c &= 100 \\
k_s &= 3.0 \\
s_s &= 1.5 \\
k_c &= 10.0 \\
s_c &= 1.5 \\
a_{\max} &= 10.0~\mathrm{m/s^2} \\
a_{\mathrm{comfort}} &= 2.0~\mathrm{m/s^2} \\
T_r &= 1.0~\mathrm{s} \\
b_{\mathrm{long}} &= 0.5~\mathrm{m} \\
b_{\mathrm{lat}} &= 1.0~\mathrm{m} \\
n &= 4 \\
v_{\mathrm{lat,min}} &= 0.1~\mathrm{m/s} \\
L_{\max} &= 15.0~\mathrm{m} \\
f_{\mathrm{lane}} &= 1.0
\end{align}

\subsection{Representative Constraints}

The commonly used scenario constraints are

\begin{align}
v &\in [0,\; 10]~\mathrm{m/s} \\
a &\in [-4,\; 2.5]~\mathrm{m/s^2} \\
\delta &\in [-0.6,\; 0.6]~\mathrm{rad} \\
\dot{\delta} &\in [-0.4,\; 0.4]~\mathrm{rad/s}.
\end{align}

\end{document}
